\section{Economics: who receives the heat value?}
\label{sec:economics}

The case break-even is an accounting measure for the combined operator and host. For one
full-load-equivalent hour, the case server consumes \qty{1}{kWh} of electricity at
CHF~0.3325/kWh and delivers \qty{0.8}{kWh} of useful heat worth
CHF~$0.3325/2.5=0.133$ per kWh. If all that heat displaces heat-pump operation, the combined
electricity cost net of heat value is $0.3325-0.8(0.133)=\text{CHF}~0.2261$ per hour.
For the four-server fleet run of the reference house, the same calculation is
$(\text{CHF}~2{,}964.84-\text{CHF}~948.75)/8{,}916.80
=\text{CHF}~0.2261$ per sold server-hour
(Appendix~\ref{app:repro}). This number excludes equipment and operation.

\paragraph{Operator and host.} In the provisional operator model, DigitalHeat
owns the hardware, sells the computing and pays or reimburses its electricity. The homeowner
receives the useful heat and retains a backup heat source. With free heat, the host's potential
avoided heat-pump expense is $Q_{\mathrm{use}}v$, while DigitalHeat's cash balance is computing
receipts minus electricity, depreciation, operations, wear and installation. The CHF~949 heat
value in the four-server run belongs to the host; it is not also operator income. If the host pays
a fee $h$ per useful kWh, $Q_{\mathrm{use}}h$ transfers from host to operator without changing
combined surplus. A positive combined surplus alone does not establish a viable offer for both.

\paragraph{Full cost under sourced demand.} With the sourced job demand of the fleet model, the
four-server reference house earns CHF~13,360 a year (CHF~8,911 from Type~1, CHF~4,449 from
Type~2). Electricity costs CHF~2,965, so the margin, which includes the heat value, is CHF~11,344. The
L4 hardware profile, normalised to a \qty{1}{kW} server at CHF~34,244 over a six-year accounting
life plus CHF~708 yearly operations, costs CHF~6,415 per server-year, or CHF~25,661 for four
servers. The net after hardware is $-\text{CHF}~14{,}317$ per year. Table~\ref{tab:breakeven}
changes one input at a time to find where the net reaches zero.

\begin{table}[ht]
  \centering
  \caption{Hardware break-even of the reference house in the fleet model (Basel 2025, L4 profile,
    sourced demand scaled by 0.1 unless stated). Net = income + heat value $-$ electricity $-$
    hardware per year.}
  \label{tab:breakeven}
  \begin{tabular}{>{\raggedright\arraybackslash}p{0.38\linewidth}>{\raggedright\arraybackslash}p{0.54\linewidth}}
    \toprule
    Change & Result \\
    \midrule
    Servers per house & 1: $+$CHF~5,252, payback 3.1 years, 18.9\% coverage, 5.9\% missed deadlines;
      2: $-$CHF~621; 3: $-$CHF~7,654; 4: $-$CHF~14,317 \\
    Job volume (all types), 4 servers & break-even at 2.8 times the scaled demand; 30\% of jobs then
      miss their deadline, and above about 3 times the margin levels off near CHF~35,000 because the
      house cannot use more heat \\
    Type~1 prices, 4 servers & break-even at 2.61 times the sourced prices (LLM inference about
      CHF~5.81 per L4 server-hour instead of CHF~2.23) \\
    Type~2 price band only & no break-even: with elasticity $-0.8$, the accepted volume falls almost
      as fast as the price rises, so revenue above the competitor price grows only with $p^{0.2}$ \\
    \bottomrule
  \end{tabular}
\end{table}

At the scaled demand, income hardly depends on the number of servers (CHF~13,300 to 14,100 for one
to four), because work is limited by demand in winter and by heat in summer, not by computing
capacity. The hardware therefore decides the result, and the single-server house is the only
configuration that covers it at this demand. All nets credit the heat value to the combined system; for the operator alone,
who gives the heat away, they are lower by that value (CHF~949 with four servers, CHF~784 with one,
so the single-server operator keeps about CHF~4,470). The price route is optimistic: Type~1
customers do not react to price in the model. The volume baseline is itself a scenario value (Section~\ref{sec:limits}).

\paragraph{Upper bound with unlimited work.} Two L4 purchase costs appear in this section: the
hardware profile's CHF~34,244 per normalised \qty{1}{kW} server, used by the fleet results above,
and a later card-price sensitivity of CHF~44,179. Neither is a quoted installable chassis. If every
heat-feasible hour were sold (the case's unlimited-work assumption, \num{22661} server-hours for four
servers), the four-server, free-heat operator would need an average realised compute price of
\begin{equation}
  p_{\mathrm{op}}^{*}
  =0.3325+\frac{4(44{,}179/6+708)}{22{,}661}
  \simeq \text{CHF}~1.76\text{/server-hour}
  \label{eq:operator-break-even}
\end{equation}
at the later cost, or about CHF~1.46 at the profile's cost, with six-year depreciation and CHF~708
yearly operations per server. The calculation excludes installation, cooling, financing, platform
fees, idle power and wear. Consumer-card deployment is excluded from the proposed architecture
pending licence clarification; the cost ledger is identified in Appendix~\ref{app:repro}.

Two token-price bases give different hourly revenue. The fleet job type uses \num{25} million tokens
per L4 server-hour and the Public AI Apertus~8B price, about CHF~2.23 per server-hour. An earlier
benchmark for another 8B model (\num{14.36} million tokens per server-hour) gives CHF~1.81 at that
competitor's list price and about CHF~0.90 at its batch price. Neither is a DigitalHeat sale: both
transfer benchmarks for a different model and GPU, and no recurring Type~1 order is documented. In
the unlimited-work scenario at the later purchase price, with all hours sold at that list price,
four servers would leave about CHF~1,172 per year before wear and installation; this is not a fleet
result. Wear is not included: the available wear estimate was derived for a different card count
and utilisation and cannot be applied per sold hour here. The CPU-only measured research kernel is
also not evidence of a full-power, billable GPU filler workload.

\paragraph{Market comparison.} The dated price comparison in Appendix~\ref{app:repro}
provides a screening test, not a sales
forecast. At the profile's L4 hardware cost, full heat credit and the unlimited-work \qty{64.7}{\percent} assumed
sold hours, it estimates CHF~0.054 per million 8B-class tokens; at \qty{35}{\percent} sold
hours (above the fleet model's 25.4\%, Section~\ref{sec:eval-house}), CHF~0.093. A listed Swiss 8B offer was CHF~0.034 per million tokens on the same
date. The estimates use benchmark throughput for a different GPU workload, compare models
with unequal quality and omit the operator/host split. Under free heat and the later L4
purchase sensitivity above, they cannot establish a profitable or price-leading offer.

\paragraph{A smaller residential candidate.} A separate exploratory variant sizes a physical
node below the home's year-round hot-water demand. Four L4 cards on an R570-class host respect
Dell's documented four-card limit; the existing host estimate plus card TDP gives about
\qty{0.4505}{kW} electrical power and, at the case's 80\% recovery assumption,
\qty{0.3604}{kW} recoverable heat. This is below the case's constant
$12/24=\qty{0.5}{kW}$ hot-water demand. At continuous operation the variant could deliver at
most \qty{3157}{kWh} useful heat, worth about CHF~420/year to the host at the case's heat-pump
comparison. These are estimated physical quantities, not measured wall power or installed heat
recovery (Appendix~\ref{app:repro}). The fixed reference house and
its abstract \qty{1}{kW} server remain unchanged.

At a retail comparison price of CHF~0.2924 per GPU-hour, four billable GPUs, the estimated
power and a six-year accounting life, that exploratory node needs approximately 5,376 paid
full-node hours per year under the 9\% repair stress scenario, before missing integration and
service costs. This is about 61\% of the year. The retail price is not an achievable DigitalHeat
net contract price; paid reservations can exist without active heat production. The proposed
experiment is therefore a measured, thermally integrated small node coupled to a minimum-volume
customer offer, not a demonstrated profitable installation. Its operator and host cash flows must
be recalculated from measured whole-node power, useful heat and net contracted receipts.
